\begin{abstract}
How do language models perform integer addition? We present a mechanistic analysis of the ones-digit addition circuit across three architectures (Gemma~2B, Phi-3~Mini, LLaMA~3.2-3B), revealing that digit identity is encoded in a \textbf{9-dimensional Fourier subspace} corresponding to the irreducible representations of $\mathbb{Z}/10\mathbb{Z}$. We establish a complete causal evidence triangle: the Fourier subspace is \emph{sufficient} (Fisher activation patching transfers 100\% of digit information), \emph{necessary} (multi-layer ablation reduces accuracy to near-chance), and \emph{specific} (matched-dimension random controls have zero effect). We discover that models \emph{progressively rotate} this subspace across layers to align with the unembedding matrix for readout, and that different architectures emphasize different components of the Chinese Remainder Theorem decomposition---Gemma emphasizes mod-2 (parity), Phi-3 emphasizes mod-5, and LLaMA uses a balanced ordinal encoding. Finally, we show that Fourier phase rotation within this subspace can \emph{steer} digit predictions with 70--88\% exact shift accuracy when combined with unembedding-informed corrections, demonstrating genuine mechanistic control over arithmetic computation.
\end{abstract}
